\section{Method}
\label{sec:method}

This section describes the protocol as implemented in the accompanying
package, version~0.2.0. Symbols in \code{typewriter} font name the
corresponding code objects.

\subsection{Data and notation}
\label{sec:data}

A \emph{panel} holds daily bars for trading sessions $t=1,\dots,T$ and assets
$i=1,\dots,N$: open $O_{t,i}$, high $H_{t,i}$, low $L_{t,i}$, close $C_{t,i}$
and volume $V_{t,i}$. The panel accepts only bars whose finality is
\code{confirmed}. For US equities these are read from the suite's shared
store through its market-data gateway; the synthetic benchmark generates bars
in the same canonical schema. Every panel carries a SHA-256 content hash over
its dates, symbols, sources and values. When built from the store, it also
carries the store manifest hash captured at read time.

\subsection{The factor language}
\label{sec:dsl}

\paragraph{Syntax.} Expressions are parsed by a hand-written recursive-descent
parser. User text is never evaluated as code. The grammar is
\begin{align*}
\mathit{expr} &::= \mathit{term}\ \big((\texttt{+}\mid\texttt{-})\ \mathit{term}\big)^{*} &
\mathit{term} &::= \mathit{unary}\ \big((\texttt{*}\mid\texttt{/})\ \mathit{unary}\big)^{*}\\
\mathit{unary} &::= \texttt{-}\,\mathit{unary} \mid \mathit{primary} &
\mathit{primary} &::= \mathit{number} \mid \mathit{name} \mid \mathit{name}\texttt{(}\mathit{expr}(\texttt{,}\,\mathit{expr})^{*}\texttt{)} \mid \texttt{(}\mathit{expr}\texttt{)}
\end{align*}
Infix operators desugar to \code{add}, \code{sub}, \code{mul} and
\code{div}. Unary minus desugars to \code{neg}, or to a negative literal when
applied to a number.

\paragraph{Terminals.} There are eight terminals:
\begin{itemize}
  \item the five bar fields \code{open}, \code{high}, \code{low},
    \code{close} and \code{volume};
  \item \code{returns} $=C_{t}/C_{t-1}-1$, undefined when $|C_{t-1}|<10^{-12}$;
  \item \code{vwap\_proxy} $=(H_t+L_t+C_t)/3$, since daily bars carry no true VWAP;
  \item \code{dollar\_volume} $=C_tV_t$.
\end{itemize}
Forward returns are \emph{not} a terminal. They exist only inside the
evaluation metrics, after the execution lag.

\paragraph{Operators.} The language has 27 operators in three classes
(Table~\ref{tab:operators}). Time-series windows $d$ must be integer
literals: $d\ge 1$ for the shifts \code{delay} and \code{delta}, and
$d\ge 2$ for rolling operators. Rolling operators use only full windows, so a
window touching a missing observation is undefined. Cross-sectional operators
read one date at a time. Non-finite intermediate values are set to missing,
and the final factor is masked to cells with a confirmed bar.

\begin{table}[t]
\centering
\small
\caption{Operators of the factor language. $x,y$ are series, $d$ an integer
window literal and $c$ a numeric literal. Signed transforms preserve the sign
of their argument.}
\label{tab:operators}
\begin{tabular}{lp{0.68\textwidth}}
\toprule
Class & Operators and semantics \\
\midrule
Element-wise (10) & \code{add}, \code{sub}, \code{mul}; \code{div} (undefined where $|y|<10^{-12}$); \code{neg}, \code{abs}, \code{sign}; signed \code{log} $=\sign(x)\log(1+|x|)$; signed \code{sqrt}; signed \code{power}$(x,c)=\sign(x)|x|^{c}$ with $c\in[-4,4]$ \\
Shifts (2) & \code{delay}$(x,d)=x_{t-d}$; \code{delta}$(x,d)=x_t-x_{t-d}$ \\
Rolling (12) & \code{ts\_mean}, \code{ts\_sum}, \code{ts\_min}, \code{ts\_max}; \code{ts\_std} (ddof $1$); \code{ts\_rank} (average-tie percentile rank of $x_t$ among the last $d$ values); \code{ts\_zscore}; \code{ts\_argmax}/\code{ts\_argmin} (sessions since the window extremum, ties to the most recent); \code{decay\_linear} (weights $d,d-1,\dots,1$ from today backwards, normalized by $d(d+1)/2$); \code{ts\_corr}, \code{ts\_cov} (on jointly observed values, ddof $1$; degenerate windows undefined) \\
Cross-sectional (3) & \code{cs\_rank} (per-date percentile rank), \code{cs\_zscore} (ddof $1$), \code{cs\_demean} \\
\bottomrule
\end{tabular}
\end{table}

\paragraph{Static validation.} A validator checks the following and returns
every violation as a machine-readable code, which is fed back to proposers:
\begin{itemize}
  \item operator names, arities and argument kinds;
  \item window and parameter ranges;
  \item depth at most 10 (a leaf has depth $1$);
  \item at most 48 nodes;
  \item windows at most 252 sessions;
  \item literals $|c|\le 10^{6}$;
  \item cumulative lookback at most 504 sessions;
  \item at least one data terminal.
\end{itemize}
The \emph{lookback} of an expression is defined recursively. It is $1$ for
\code{returns}, $0$ for the other terminals and for literals, and for an
operator $f$ with window $d$ applied to arguments $x_j$ it is
\[
\lambda\big(f(x_1,\dots;d)\big)=\max_j \lambda(x_j)+\omega_f(d),
\]
where $\omega_f(d)=d-1$ for rolling operators, $\omega_f(d)=d$ for shifts and
$\omega_f=0$ otherwise.

\paragraph{Causality.} Every operator reads the current row and a bounded
number of earlier rows of its inputs, or a single date across assets.
Structural induction therefore gives the following property: for any valid
expression $e$, the value $e_{t,i}$ is a function of bars dated at or before
$t$ only. The test suite checks this property empirically. It perturbs all
bars after a date $\tau$ and verifies that no factor value at or before
$\tau$ changes, for every operator, every terminal, nested composites and the
reference library.

\paragraph{Canonical form, hashing and complexity.} The canonical form sorts
the two series arguments of the commutative operators (\code{add},
\code{mul}, \code{ts\_corr}, \code{ts\_cov}) by their canonical strings, and
writes literals in non-window slots as floats. The structural hash is the
SHA-256 of a namespaced canonical string. It is the deduplication key and
the evaluator's cache key. Syntax alone cannot identify equivalent
hypotheses: $(a+b)+c$ and $a+(b+c)$ hash differently, and so do
\code{volume}, \code{log(volume)} and \code{cs\_rank(volume)}, whose
per-date ranks are identical. Equivalence for multiple testing is therefore
decided on behaviour (Section~\ref{sec:protocol}). The complexity of an
expression is
$\kappa(e)=n_{\mathrm{op}}+n_{\mathrm{term}}+\tfrac12 n_{\mathrm{lit}}$.

\subsection{Evaluation timing and metrics}
\label{sec:metrics}

\paragraph{Execution lag.} A signal $f_{t,i}$ observed at the close of $t$ is
traded at the close of $t+\ell$ and held for $h$ sessions. The forward
return is
\[
y^{(h)}_{t,i}=\frac{C_{t+\ell+h,i}}{C_{t+\ell,i}}-1,\qquad \ell\ge 1,
\]
and $\ell\ge1$ is enforced in code. All metric series are indexed by the
signal date $t$. We use $\ell=h=1$ throughout.

\paragraph{Windows and embargo.} The calendar is split chronologically by
session counts into formation, validation and test windows. The default
fractions are $0.6/0.2/0.2$ of the sessions left after two embargo gaps. The
windows are separated by $E\ge\ell+h_{\max}$ sessions that belong to no
window. Within a window $W$ with last session position $p_W$, only signal
positions with $p+\ell+h\le p_W$ are scored, so no forward return ends
outside its window. Because factors are causal, each phase computes a factor
on the panel it is allowed to see and then restricts the result to the
window's scored dates. The value on any date is the same whichever panel was
used, so there is no recomputation effect at window boundaries.

\paragraph{Information coefficient.} The per-date rank IC is the Pearson
correlation of average-tie ranks of $f_{t,\cdot}$ and $y_{t,\cdot}$ over
assets where both are observed. It is undefined on dates with fewer than
$10$ such assets or no rank variation. For the IC series on a window's
scored dates:
\begin{itemize}
  \item $\overline{\IC}$ is the mean;
  \item $\ICIR=\overline{\IC}/\mathrm{sd}(\IC)$ is per period, not
    annualized;
  \item $t_{\mathrm{NW}}=\overline{\IC}/\widehat{\mathrm{se}}_{\mathrm{NW}}$
    uses the Newey--West standard error with a Bartlett
    kernel~\citep{newey1987hac} and $q=\max\{h-1,\lfloor 4(n/100)^{2/9}\rfloor\}$
    lags, where $n$ is the number of IC dates.
\end{itemize}
p-values refer $t_{\mathrm{NW}}$ to a $t$ distribution with $n-1$ degrees of
freedom, a conservative small-sample convention.

\paragraph{Portfolio metrics.} On each date, assets are split into five
equal-count buckets by signal. The long-short portfolio holds $+1$ spread
equally over the top bucket and $-1$ over the bottom bucket. With $h>1$, the
weights held are the average of the last $h$ formations. Weights held on
signal date $t$ earn the close-to-close return from $t+\ell$ to $t+\ell+1$.
Trading costs are $10$ basis points per unit of $\sum_i|w_{t,i}-w_{t-1,i}|$.
Sharpe ratios are annualized with $\sqrt{252}$. Top-bucket turnover is the
share of today's top-bucket names that were not in yesterday's top bucket.

\subsection{Search, selection and the sealed hold-out}
\label{sec:protocol}

\paragraph{Phase 1: search.} The search loop evaluates proposals on a panel
truncated at the end of the formation window, so every statistic a proposer
can receive is computed from formation data. The harness can also be given a
loader instead of a panel; each phase then reads only its own data, and the
test window is loaded only after the factor set is frozen. The store path of
the command-line interface works this way. Each round proceeds as follows.
\begin{enumerate}
  \item The proposer receives a context (described below) and returns a batch
    of $b=20$ expressions, or the remaining budget if smaller.
  \item Each expression is validated. An invalid expression becomes a
    \emph{trial} with its error codes. If its whitespace- and
    case-normalized text repeats an earlier invalid text, it is logged as a
    duplicate instead.
  \item A valid expression is canonicalized. If its hash was seen before, it
    is logged as a duplicate and consumes no budget. Otherwise it is
    evaluated on the formation window and becomes a trial. A numerical
    failure is also a trial, with status \code{error}.
  \item An evaluated trial whose IC is defined on fewer than
    $\max(\SynMinIcDates, \lceil \SynMinIcFraction\, n_{\mathrm{scorable}}\rceil)$
    formation dates is \emph{degenerate}: it is a trial with $p=1$, and the
    proposer is told so. Without this rule, and with a normal reference, a
    factor with three defined IC dates can obtain an arbitrarily small
    p-value.
\end{enumerate}
The loop stops for any of these reasons:
\begin{itemize}
  \item the budget of $B$ trials is reached;
  \item $\max(100, 10\lceil B/b\rceil)$ rounds have run (100 at $B=200$,
    $b=20$), so that larger budgets can be used in full;
  \item 5 consecutive rounds produced no new trial;
  \item 3 consecutive proposer failures occurred (for example, refusals).
\end{itemize}
Proposer failures are logged and do not count as trials. Only a run that
reaches its budget is \emph{complete}. A run that ends with repeated proposer
failures, or with no trial at all, is \emph{aborted}: it is logged but
neither committed nor revealed. Misconfiguration (a rejected request, key,
permission or model identifier) is not a proposer failure; it stops the run
with an error, so it cannot pass for a null result.

\paragraph{Information barrier.} The barrier is an interface: it limits what
the harness passes to the proposer. It is not a sandbox; in-process code
could read anything in memory. For the LLM arm, only the rendered prompt
leaves the process. The context type is frozen and slotted. Its feedback
fields accept only two exact types, and subclasses are rejected:
\begin{itemize}
  \item evaluated feedback: canonical text, formation IC, ICIR,
    $t_{\mathrm{NW}}$, top-bucket turnover and complexity. The only factory
    that builds it from a metrics report refuses any window not named
    \code{formation};
  \item rejected feedback: the proposal's text, its error codes and a
    message, truncated to 400 characters. For an invalid proposal the message
    is the validator's; a failed or degenerate trial gets a fixed text (for a
    failure, only the exception type) without the date counts that the
    ledger records.
\end{itemize}
Each round's context contains:
\begin{itemize}
  \item the grammar card and the active limits;
  \item the 8 trials with the highest $|\ICIR|$;
  \item the 4 with the lowest, once more than 8 exist;
  \item the previous round's evaluations (the LLM prompt shows at most 10);
  \item the last 8 rejections;
  \item the trial count and the remaining budget.
\end{itemize}
It contains no dates, asset identifiers, date counts, or validation or test
quantities. The formation statistics still imply the approximate formation
sample size, because $t_{\mathrm{NW}}/\ICIR$ grows roughly like $\sqrt{n}$. An end-to-end test perturbs all post-formation bars and confirms
that every context a proposer receives is unchanged.

\paragraph{Phase 2: selection.} Selection loads a panel truncated at the end
of the validation window.
\begin{enumerate}
  \item \emph{Behavioural classes.} Two evaluated trials whose per-date
    cross-sectional ranks on the formation IC dates agree up to sign have
    identical rank-IC series against any forward return, so they test one
    hypothesis. Each such class is represented by its first trial. Counting
    the copies separately would make Benjamini--Hochberg \emph{less} strict,
    because a step-up rule rewards clusters of equal p-values.
  \item \emph{Formation screen.} The Benjamini--Hochberg step-up
    procedure~\citep{benjamini1995fdr} at $q=\SynFdrAlpha$ is applied to the
    two-sided formation p-values of the class representatives, with family
    size $m=\max(B,n_{\mathrm{trials}})-n_{\mathrm{dup}}$: every budgeted
    slot that did not repeat a tested hypothesis counts, degenerate, invalid
    and failed trials and unused budget as $p=1$, so a run that stops early
    faces the family of a full one. A feedback-driven proposer chooses its
    next expressions after seeing formation results, so these p-values are
    not valid for false-discovery control. The screen is a filter only.
    The implementation also offers
    Benjamini--Yekutieli~\citep{benjamini2001fdr} and Holm's step-down
    procedure. \todo{add a citation for Holm's procedure once its metadata
    has been verified.}
  \item \emph{Orientation and confirmation.} Each survivor $j$ is multiplied
    by $s_j=-1$ if its formation IC is negative and $s_j=+1$ otherwise, and
    scored on the validation window, which no proposer has seen.
    Benjamini--Hochberg at $q=\SynConfirmAlpha$ is applied to the one-sided
    validation p-values of all survivors, with $m$ equal to the number of
    survivors; a survivor that is degenerate on validation gets $p=1$. This
    is the step whose error control carries the claim: given the candidate
    set, which depends on formation data only, the validation p-values are
    valid.
  \item \emph{Decorrelation.} Confirmed candidates are taken in descending
    order of validation ICIR, with ties broken by hash. A candidate is kept
    if the mean over validation dates of its cross-sectional Spearman
    correlation with every kept factor is below $0.7$ in absolute value. At
    most $k=5$ factors are kept.
\end{enumerate}
For each survivor we also record the deflated Sharpe
ratio~\citep{bailey2014deflated} of its formation ICIR as a
\emph{diagnostic}; it is not used for selection. Survivors are selected on
$|\ICIR|$, a two-sided choice: under a symmetric null the maximum of $N$
absolute values behaves like the maximum of $2N$ signed ones, so the null
threshold uses $2N$ trials,
\[
\widehat{SR}_0=\sqrt{V}\Big[(1-\gamma)\,\Phi^{-1}\!\big(1-\tfrac1{2N}\big)+\gamma\,\Phi^{-1}\!\big(1-\tfrac{1}{2Ne}\big)\Big],
\]
where $V$ is the sample variance of the class representatives' formation
ICIRs, $N$ the screen's family size and $\gamma$ the Euler--Mascheroni
constant. The diagnostic is then
\[
\mathrm{DSR}=\Phi\!\left(\frac{(\widehat{SR}-\widehat{SR}_0)\sqrt{n-1}}{\sqrt{1-g_3\widehat{SR}+\tfrac{g_4-1}{4}\widehat{SR}^2}}\right),
\]
with $\widehat{SR}=|\ICIR|$, and $g_3$, $g_4$ the skewness (sign-adjusted by
$s_j$) and non-excess kurtosis of the IC series. Treating trials as
independent makes this diagnostic approximate.

\paragraph{Phase 3: commit and reveal.} The selected set $F$ is an ordered
tuple of canonical expressions, hashes and orientations. The seal
identifier is $\sigma=H(\text{terms})$, where $H$ is SHA-256 over canonical
JSON and the terms are the test window, the panel content hash, the metric
configuration and the language limits. The commitment
$c=H(\sigma,\text{terms},F,H(F))$ is appended to the ledger before any test
metric exists. The one-commit, one-reveal rule is keyed on a \emph{guard}
$g=H(\text{data hash},\text{test-window dates})$ only, so changing the metric
settings or the limits yields a new $\sigma$ but not a second reveal of the
same window. The ledger is re-read on every commit and reveal, so several
instances, or several writers of one ledger file, share the rule. A commit is
refused if the guard already holds a commitment or a reveal, or if the data
hash recomputed from the panel's values differs from the one it was built
with. A reveal recomputes the commitment for the set it is given and refuses
in four cases:
\begin{itemize}
  \item nothing has been committed;
  \item a reveal for $g$ already exists in the ledger;
  \item the set or the seal terms differ from the committed ones;
  \item an external authority refuses it (below).
\end{itemize}
The last three refusals are written to the ledger. On a successful reveal,
test metrics are computed for each factor and for the composite. The
composite is the equal-weight, missing-aware mean of per-date z-scored
oriented signals.

\paragraph{Study registry.} Across runs, a project-wide register fixes each
real-data \emph{study} (symbols, dates, split and a pre-registered maximum
number of reveals) before anything is evaluated, counts every exploratory
evaluation outside a search, records commitments and caps reveals. On store
data the search only commits; revealing is a separate command that requires
the commitment hash, which is meant to be published outside the repository
first, and the registry's permission. The registry, like the ledger, is
self-attested: it makes repeated reveals countable, not impossible.

\paragraph{Trial ledger.} The ledger is an append-only JSON Lines file. Record
$k$ stores
\[
h_k=H(\text{version},k,\text{kind}_k,\tau_k,\text{payload}_k,h_{k-1}),\qquad h_{-1}=0^{64},
\]
where $\tau_k$ is an optional timestamp. Canonical JSON uses sorted keys and
rejects NaN. The ledger records the run configuration, every trial,
duplicate, proposer error and round, the selection with all candidates and
decisions, the commitment, the reveal and the run end. Any edit, insertion,
deletion or reordering breaks the chain. Truncating the tail is detected by
comparing against a separately stored head hash. These checks detect
accidental or naive edits only: the head is written by the same process,
into the same directory, so anyone with write access can drop records and
rebuild the chain. A head or commitment hash proves something to a third
party only once it has been published elsewhere before the reveal.

\subsection{Proposers}
\label{sec:proposers}

\paragraph{Random grammar (A).} This baseline samples typed trees i.i.d. and
ignores all feedback.
\begin{itemize}
  \item Maximum depth is 4.
  \item A series slot below the root becomes a leaf with probability $0.35$.
  \item The second argument of an arithmetic operator becomes a literal from
    $\{0.5,1,2\}$ with probability $0.15$.
  \item Windows are drawn from $\{1,2,3,5,10,20,40,60\}$ (shifts may use
    $1$; rolling operators need at least $2$) and exponents from
    $\{0.5,2\}$. The sampling probability of any tree is computable, which
    lets a test confirm that every planted benchmark signal is in the
    sampler's support.
  \item Sampling repeats until the tree is valid. It is then redrawn, up to
    50 times, if it repeats one of the proposer's earlier samples.
\end{itemize}

\paragraph{Genetic programming (B).} This baseline follows the classic
setup~\citep{koza1992genetic}. It keeps an archive of evaluated trees with
fitness $|\ICIR_{\mathrm{formation}}|-0.002\,\kappa$; the sign is free
because selection orients factors. Each round:
\begin{itemize}
  \item the 30 fittest archive members form the population;
  \item children are bred by tournament selection (size 3) and one of three
    operations: subtree crossover ($p=0.5$), subtree mutation ($0.3$, with
    new subtrees of depth at most 3) or point mutation ($0.2$), which swaps
    an operator for another of the same signature, a terminal, or a window;
  \item a tenth of each batch is random immigrants.
\end{itemize}
Children must be valid and new. These operator settings were fixed before
any experiment and were not tuned on results. The window menu that both
baselines share was extended by the value $1$ in the revision disclosed in
Section~\ref{sec:limitations}, after the first harness runs had been seen,
so that the library-family planted signal's one-session change became
reachable; before, it was $\{2,3,5,10,20,40,60\}$.

\paragraph{LLM-guided (D).} The prompts are two versioned templates, a system
prompt and a user prompt (version~2), whose SHA-256 hashes are recorded with
every call and every proposal. Version~2 dropped the example mechanism labels
of version~1, which named the stories behind two planted benchmark signals.
The system prompt states four things:
\begin{itemize}
  \item the data are anonymized;
  \item the model should not try to infer the market or period;
  \item signs are free;
  \item short, economically motivated expressions are preferred.
\end{itemize}
The user prompt carries the grammar card, the formation feedback blocks
(best, worst and previous round), the recent rejections and the remaining
budget. The response must match a JSON schema with exactly the fields
\code{expression}, \code{rationale} and \code{economic\_mechanism} per
proposal, and no additional properties.

Before sending, the harness-written text of every prompt (templates and
grammar card) is checked for:
\begin{itemize}
  \item ISO-style dates;
  \item month names (full names in any case; ``March'', ``May'' and
    abbreviations when capitalized);
  \item four-digit year-like numbers, excluding the integer counters;
  \item any whole-word occurrence of the panel's symbols.
\end{itemize}
A violation aborts the run. Model-written text echoed back in the rejection
block is checked for the same patterns (except years) and redacted rather
than aborting the run, since it is the model's own output.

The Claude backend sends the following:
\begin{itemize}
  \item the model identifier (default \code{claude-opus-5}) and a maximum of
    16{,}000 output tokens;
  \item adaptive thinking;
  \item an effort level (default \code{high}) and the JSON-schema output
    format.
\end{itemize}
It sends no sampling parameters and no prefilled assistant turn. Responses
are handled as follows:
\begin{itemize}
  \item a refusal stop reason is checked before content is read;
  \item rate limits, overload and server errors, and connection failures are
    retried with bounded exponential backoff;
  \item a rejected request, key, permission or model identifier (HTTP 400,
    401, 403, 404, 413 or 422) stops the run;
  \item truncation at the token limit, invalid JSON and failures that
    survived the retries raise recoverable proposer errors;
  \item paused turns are resumed up to three times.
\end{itemize}
The served model identifier, stop reason, token usage and request
identifier are recorded for every call. Server-side model fallback can
silently change the model under test, so it is an explicit option that is
off in all experiments. A recording layer appends every request and its
outcome to a write-only JSON Lines file, in call order and failures included.
It is keyed by the hash of the request, including a replicate tag and a tag
for the benchmark cell, and of the backend identity. A replay backend
answers the $n$-th occurrence of a request with the $n$-th recorded outcome
and treats anything unrecorded as fatal, so a replayed run follows the live
run's path exactly; a test checks this for a live run that met a transient
failure.

\subsection{Contamination diagnostics}
\label{sec:contamination}

\paragraph{Knowledge-cutoff split.} Given a per-date IC series and a cutoff
session $c$ from the model provider's documentation, the split works as
follows:
\begin{itemize}
  \item \emph{pre} uses signal positions with $p+\ell+h<c$;
  \item \emph{post} uses positions with $p\ge c$;
  \item signal dates whose forward return straddles the cutoff are dropped.
\end{itemize}
The difference in mean IC is standardized by
$\sqrt{\widehat{\mathrm{se}}_{\mathrm{pre}}^2+\widehat{\mathrm{se}}_{\mathrm{post}}^2}$,
with Newey--West standard errors. Section~\ref{sec:setup} uses this split in
a difference-in-differences design against a baseline. No cutoff date is
assumed anywhere in the code.

\paragraph{Novelty.} Rediscovery is measured on values. The
\emph{behavioural novelty} of a factor $f$ against the reference library
$\mathcal{L}$ is
\[
1-\max_{\ell\in\mathcal{L}}\Big|\frac{1}{|T|}\sum_{t\in T}\rho_S\big(f_{t,\cdot},\ell_{t,\cdot}\big)\Big|,
\]
the largest mean cross-sectional Spearman correlation with any library
entry over the formation IC dates $T$. A re-spelling of a classic factor,
such as \code{-ts\_sum(returns, 5)} for five-day reversal, scores near
$0$. \emph{Structural novelty}, reported as a secondary description of
syntax, is $1-\max_{\ell}|S(e)\cap S(\ell)|/|S(e)\cup S(\ell)|$ with $S(e)$ the
set of structural hashes of $e$'s operator and terminal subtrees; it cannot
tell a re-spelled classic factor from new structure. The library has 18
entries covering momentum, reversal, volatility, liquidity and volume-price
interaction. Three are translations of published formulas of
\citet{kakushadze2016alphas} (Alphas \#2, \#6 and \#13); four more are
the project's own variants in that style.
